\SourceLeaf{255}%
woraus leicht folgt, dass die Gleichungen (5) und (6) für unseren Fall vollständig ausreichen, sofern das, was sie unbestimmt lassen, so gewählt wird, dass \(r\) und \(\varphi\) für den Anfangspunkt und für Punkte in unendlich kleinem Abstand von ihm die richtigen Werte annehmen.

Was die Integration der Gleichungen (5) und (6) selbst betrifft, so lässt sie sich bekanntlich auf die Integration gewöhnlicher Differentialgleichungen zurückführen. Diese werden jedoch meist so verwickelt, dass wenig gewonnen ist. Die Entwicklung in Reihen bereitet dagegen keine Schwierigkeiten. Die Reihen genügen für praktische Zwecke bei kleinen Teilen der Fläche vollkommen, und die angegebenen Formeln bieten damit eine ergiebige Grundlage zur Lösung vieler sehr wichtiger Probleme. Hier werden wir jedoch nur ein einziges Beispiel ausarbeiten, um die Art der Methode zu zeigen.
\Article{23}
Betrachten wir den Fall, dass alle Linien, auf denen \(p\) konstant ist, kürzeste Linien sind, die die Linie mit \(\varphi=0\) rechtwinklig schneiden. Diese können wir als Abszissenlinie auffassen. Es sei \(A\) der Punkt, für den \(r=0\) gilt, \(D\) ein beliebiger Punkt der Abszissenlinie und \(AD=p\). Ferner sei \(B\) ein beliebiger Punkt auf der kürzesten Linie, die senkrecht auf \(AD\) steht, und zwar in \(D\), und es sei \(BD=q\). Dann lässt sich \(p\) als Abszisse und \(q\) als Ordinate von \(B\) auffassen. Auf dem Zweig der Abszissenlinie, dem \(\varphi=0\) entspricht, rechnen wir die Abszissen positiv. Dabei betrachten wir \(r\) stets als positive Größe. Die Ordinaten nehmen wir auf der Seite positiv, auf der \(\varphi\) zwischen 0 und \(180^0\) gemessen wird.

Nach dem Satz aus Artikel 16 gilt \(\omega=90^0\), \(F=0\). Außerdem gilt \(G=1\), und wir setzen \(\BareRoot{E}=n\). Somit ist \(n\) eine Funktion von \(p,q\). Für \(q=0\) muss sie \(=1\) sein. Wenden wir die in Artikel 18 angegebene Formel auf unseren Fall an, so zeigt sich, dass auf \textit{jeder} kürzesten Linie \(\mathrm{d}\theta=-\frac{\mathrm{d}n}{\mathrm{d}q}\mathord{.}\mathrm{d}p\) gelten muss. Dabei bezeichnet \(\theta\) den Winkel zwischen einem Element dieser Linie und dem Element der Linie, auf der \(q\) konstant ist. Da die Abszissenlinie selbst eine kürzeste Linie ist und auf ihr überall \(\theta=0\) gilt, muss für \(q=0\) überall \(\frac{\mathrm{d}n}{\mathrm{d}q}=0\) gelten. Entwickeln wir also \(n\) in eine Reihe nach aufsteigenden Potenzen von \(q\), so muss diese die folgende Form haben
\[
n=1+fqq+gq^3+hq^4+\text{etc.}
\]
wobei \(f,g,h\) usw. Funktionen von \(p\) sind. Wir setzen
